% ============================================================
%  Chapter 7 — Search Engines
%  The Secret Life of the Internet
% ============================================================

\chapter{Search Engines}

\chapteropening{%
  There are billions of pages on the internet.
  How do you find exactly the one you need?
  Scout Search has a very, very big index!
}
\chapterbadge{Detective Time with Scout Search}

% --- Opening story ---
\section*{The World's Biggest Library}

Imagine a library with one billion books.
Walking along the shelves looking for one title could take years.

Sam had a question: ``How do rainbows form?''
But there were billions of pages on the internet.
Where to start?

A figure appeared, wearing a long coat and holding a giant flashlight and a stack of index cards.
That was Scout Search.

\character{Scout Search}{Don't worry! 
  I've already visited every page on the internet,
  read them all, and made a huge organised list.
  Tell me your question and I'll find the best answers in a flash!}

\sectionrule

% --- What is a search engine? ---
\section*{What Is a Search Engine?}

A \term{search engine} is a service that helps you find information on the internet.

\begin{funfact}
Google's search index contains hundreds of billions of web pages and is over 100 million gigabytes in size!
\end{funfact}

You type a question or some keywords into a search box.
The search engine looks through its enormous index and shows you a list of pages
that are likely to help.

\begin{picturethis}
  Imagine a super-organised librarian who has read every single book in the library.
  She has written a card for every word in every book,
  saying which books contain that word and on which page.
  When you ask for help, she checks the cards and brings you the most useful books first.
  That librarian is a search engine.
  The cards are the \textbf{index}.
\end{picturethis}

\sectionrule

% --- How does a search engine work? ---
\section*{How Does a Search Engine Work?}

A search engine has three main jobs:

\subsection*{1. Crawling}

Special programmes called \term{crawlers} (or ``spiders'' or ``bots'') roam the internet
constantly, visiting web pages.
When a crawler finds a page, it follows all the links on that page to find more pages,
then follows the links on those pages — and so on, endlessly.

\begin{picturethis}
  Imagine a spider walking across a huge web.
  Every point on the web connects to more threads.
  The spider visits every point, reads what is there, and notes it down.
  That is what a web crawler does.
\end{picturethis}

\subsection*{2. Indexing}

After crawling, the search engine builds an \term{index} —
a massive organised list of words, phrases, and topics,
linked to all the pages where they appear.

This index is so large it would fill millions of filing cabinets,
but computers can search through it in less than a second.

\subsection*{3. Ranking}

When you type a search query, the search engine does not just list every matching page.
It \term{ranks} them — it tries to show you the most useful, trustworthy, and relevant pages first.

Ranking is decided by many signals, including:
\begin{itemize}
  \item How many other websites link to a page (if lots of people point to it, it is probably useful).
  \item How well the words on the page match your question.
  \item How new or up-to-date the page is.
  \item How easy the page is to read.
\end{itemize}

\begin{diagram}[From crawling pages to ranking the best results.]
  \resizebox{\linewidth}{!}{%
  \begin{tikzpicture}[>=Stealth, node distance=1.3cm]
    \node[draw=InternetBlue, rounded corners, fill=SkyBlue!25, minimum width=2.5cm, minimum height=1cm] (web) {Web Pages};
    \node[draw=WarmOrange!80!black, rounded corners, fill=WarmOrange!20, minimum width=2.4cm, minimum height=1cm, right=of web] (crawl) {Crawling \faSpider};
    \node[draw=LeafGreen!80!black, rounded corners, fill=LeafGreen!20, minimum width=2.2cm, minimum height=1cm, right=of crawl] (index) {Index};
    \node[draw=SoftPurple!80!black, rounded corners, fill=SoftPurple!20, minimum width=2.4cm, minimum height=1cm, right=of index] (rank) {Ranking \faTrophy};
    \node[draw=CoralRed!80!black, rounded corners, fill=CoralRed!16, minimum width=2.8cm, minimum height=1cm, right=of rank] (results) {Search Results};

    \draw[thick,->,InternetBlue] (web) -- (crawl);
    \draw[thick,->,InternetBlue] (crawl) -- (index);
    \draw[thick,->,InternetBlue] (index) -- (rank);
    \draw[thick,->,InternetBlue] (rank) -- (results);
  \end{tikzpicture}%
  }
\end{diagram}

\sectionrule

% --- Search tips ---
\section*{Searching Smarter}

\character{Scout Search}{Not all search words are equal!
  The better your search words, the better the results.
  Think of it like describing something clearly to a new friend.}

\begin{quickmap}
  \textbf{Good vs better search words:}
  \begin{center}
    \begin{tabularx}{\linewidth}{>{\bfseries}l Y}
      Vague search & Better search \\
      ``rainbow''  & ``how do rainbows form'' \\
      ``animal''   & ``largest animal in the ocean'' \\
      ``fix''      & ``how to fix a bicycle puncture'' \\
      ``food''     & ``healthy snack ideas for kids'' \\
    \end{tabularx}
  \end{center}
\end{quickmap}

Specific words help the search engine understand exactly what you want.

\sectionrule

% --- Being a critical reader ---
\section*{Are All Results True?}

Search engines are very good at finding pages — but not every page on the internet is accurate.

\begin{itemize}
  \item Look for well-known, trusted sources: encyclopaedias, museums, universities.
  \item If something seems surprising, check it in more than one place.
  \item Ask a trusted adult if you are not sure whether a result is reliable.
\end{itemize}

\begin{bigidea}
  Search engines are incredibly helpful, but you still need to think!
  Being a good searcher also means being a \textbf{careful reader}.
\end{bigidea}

\sectionrule

\begin{newwords}
  \begin{description}[labelwidth=3.2cm, leftmargin=3.4cm]
    \item[\term{Search engine}] A service that helps you find information on the internet by searching an index.
    \item[\term{Crawler}]       A programme that visits web pages automatically to gather information for the index.
    \item[\term{Index}]         The huge organised list of words and pages that a search engine builds.
    \item[\term{Ranking}]       The process of sorting search results to show the most useful ones first.
    \item[\term{Query}]         The words or question you type into a search box.
  \end{description}
\end{newwords}

\sectionrule

\begin{tryit}
  \textbf{Write the best search words!}

  For each question below, write the search words you would use:

  \begin{enumerate}
    \item You want to know what pandas eat.
    \item You want to find a recipe for banana bread.
    \item You want to learn the name of the tallest mountain.
    \item You want to know how a rainbow forms.
    \item You want to find out when the next solar eclipse will happen.
  \end{enumerate}

  \emph{Compare your answers with a friend.
  Whose search words do you think Scout Search would find most useful?}
\end{tryit}

\sectionrule

\begin{bigidea}
  A \textbf{search engine} crawls the internet, builds a giant index, and ranks results
  so you can find the most useful pages in seconds.
  Good search words and careful reading are your best tools.
\end{bigidea}

\character{Scout Search}{Now you know how to search like a pro!
  Next up — what IS the cloud?  Is it really made of water vapour?
  Spoiler: not quite.  Cloud Pal will explain!}
